% Fresh verified AE export; paper profile.
% Recorded platform: Linux / x86_64; structured backend: avx512vl. Hardware equivalence to the manuscript is not asserted.
% Per-setting host pairing; distributed collection. See canonical.json placement.cases.
% Recorded platform: Linux / x86_64; structured backend: avx512vl. Hardware equivalence to the manuscript is not asserted.
% Legacy xeon keys do not identify the measurement hardware.
% Linux-result subset (legacy xeon keys): unavailable PMU/M4 macros are intentionally undefined; no archived values.
% Shared performance numbers for 00-abstract.tex, 01-intro-han.tex, and 04-evaluation.tex.
% Automation: update the SECOND argument of \DefinePerfValue{key}{value}.
% One declaration per line; values are bare decimal numbers with original precision.
% Units, percent signs, and speedup symbols stay in the manuscript.
% Prefer loading this file in the main preamble, before any list of figures/tables.
% Section-level loaders also support inputting the three sections independently.
% Key: platform.experiment.[d<distance>.p<percent with p as decimal point>].variant.metric
% Example: p0p05 means physical error rate 0.05%, not probability 0.05.
% Xeon timing/speedup values use the completed Batch3 matrix; cache-miss values use independent PMU runs.
% Existing values are formatting placeholders until migration succeeds; see the generated audit JSON.
% Workload d/p, setting count, real-time target, hardware/software versions, and theoretical bounds stay literal.
% The predecoder-strategy table and narrative load separate generated predecoder_numbers.tex.
% Load this shared API file before that supplemental macro file.
% d33/p0p1 etc. inside keys identify fixed workloads; they are not adjustable d/p parameters.
% L2/L3 miss values are percentages from the completed independent PMU runs.
% Figures are external PDFs: editing this file does not change their embedded numbers.

% Idempotent loading when multiple manuscript sections are input.
\ifdefined\PerfValue\expandafter\endinput\fi

\newcommand{\DefinePerfValue}[2]{%
  \expandafter\newcommand\csname BBPerfValue@#1\endcsname{#2}%
}
\DeclareRobustCommand{\PerfValue}[1]{%
  \ifcsname BBPerfValue@#1\endcsname%
    \csname BBPerfValue@#1\endcsname%
  \else%
    \PackageError{performance-macros}{Undefined performance key '#1'}%
      {Declare the key using \string\DefinePerfValue\space before use.}%
  \fi%
}

% Abstract/introduction summaries; legacy key namespace retained for manuscript compatibility.
% Xeon workload: d=9/13/17/21/25/29/33; p=0.05%, 0.1%, 0.3%; all 21 settings.
% Single-core=Predecode; eight-core=Spatial. The d33/p0.1% example uses Spatial; round time=block time/33.
\DefinePerfValue{legacy.single-core.vs-pymatching.geomean-speedup}{4.20}
\DefinePerfValue{legacy.single-core.vs-pymatching.max-speedup}{5.58}
\DefinePerfValue{legacy.eight-core.vs-pymatching.max-speedup}{9.74}
\DefinePerfValue{legacy.d33.p0p1.beyond-blossom.mean-block-time-us}{21.79}
\DefinePerfValue{legacy.d33.p0p1.beyond-blossom.mean-round-time-us}{0.66}

% Xeon decoder-comparison figure
% Best configuration per setting; maxima over code distances at each physical error rate.
% Use the same completed Xeon matrix and Batch3 statistic as the ablation below.
\DefinePerfValue{xeon.decoder-comparison.p0p05.best.vs-pymatching.max-speedup}{8.96}
\DefinePerfValue{xeon.decoder-comparison.p0p1.best.vs-pymatching.max-speedup}{9.74}
\DefinePerfValue{xeon.decoder-comparison.p0p3.best.vs-pymatching.max-speedup}{7.22}

% Xeon incremental ablation
% Keys identify workload, configuration, and metric. Relative speedups name their denominator.
\DefinePerfValue{xeon.ablation.d33.p0p3.pymatching.mean-time-us}{947.66}
\DefinePerfValue{xeon.ablation.d33.p0p3.structured-simd.mean-time-us}{651.93}
\DefinePerfValue{xeon.ablation.d33.p0p3.structured-simd.vs-pymatching.speedup}{1.45}

% PMU percentages: 100 * miss requests / read requests; separate million-shot, LTO OFF runs.
\DefinePerfValue{xeon.ablation.d33.p0p3.pymatching.L3miss}{0.116}
\DefinePerfValue{xeon.ablation.d33.p0p3.structured-simd.L3miss}{0.019}
\DefinePerfValue{xeon.ablation.d33.p0p3.pymatching.L2miss}{60.157}
\DefinePerfValue{xeon.ablation.d33.p0p3.structured-simd.L2miss}{27.762}
\DefinePerfValue{xeon.ablation.d9.p0p05.pymatching.mean-time-us}{1.205}
\DefinePerfValue{xeon.ablation.d9.p0p05.structured-simd.mean-time-us}{1.168}
\DefinePerfValue{xeon.ablation.predecoding.vs-structured-simd.min-speedup}{2.40}
\DefinePerfValue{xeon.ablation.predecoding.vs-structured-simd.max-speedup}{4.83}
\DefinePerfValue{xeon.ablation.d13.p0p1.predecoding.mean-time-us}{1.847}
\DefinePerfValue{xeon.ablation.d13.p0p1.component-parallel.mean-time-us}{2.707}
\DefinePerfValue{xeon.ablation.d33.p0p1.predecoding.mean-time-us}{46.118}
\DefinePerfValue{xeon.ablation.d33.p0p1.component-parallel.mean-time-us}{32.178}
\DefinePerfValue{xeon.ablation.d33.p0p1.component-parallel.vs-predecoding.speedup}{1.43}
\DefinePerfValue{xeon.ablation.d33.p0p3.tail-reset.mean-time-us}{158.898}
\DefinePerfValue{xeon.ablation.d33.p0p3.parallel-predecoding.mean-time-us}{131.254}
\DefinePerfValue{xeon.ablation.d33.p0p3.parallel-predecoding.vs-pymatching.speedup}{7.22}
\DefinePerfValue{xeon.ablation.tail-reset.num-improved-settings}{11}

% Apple M4 Pro comparison and real-time coverage
% Equal-valued results for different error rates/decoders remain independent keys.
